\subsection{Finite prefixes and exact cyclic momentum}
\label{sec:ground_state_finite_prefix}

\begin{lemma}[Periodic representation of an exactly invariant finite vector]
    \label{lem:periodic_finite_cyclic_representation}
    \lean{MPSTensor.exists_mpv_eq_of_cyclic_invariant,
        MPSTensor.exists_mpvState_eq_of_cyclic_invariant,
        MPSTensor.mpvState_cyclic_invariant,
        MPSTensor.exists_mpvState_eq_iff_cyclic_invariant}
    \leanok
    \uses{thm:seqanc_cut_rank_attained, thm:seqanc_cut_rank_bounds,
        thm:pgvwc07_site_independent_matrices}
    Let $N>0$ and let $\psi\ne0$ be a finite vector satisfying
    $\psi(\sigma\circ\tau_s)=\psi(\sigma)$ for every cyclic translation
    $\tau_s$. There is a periodic trace-MPS representation of $\psi$
    with matrix size $N d^{\lfloor N/2\rfloor}$. Conversely every
    periodic trace-MPS vector has this exact cyclic invariance.
    Thus a nonzero positive-length vector is representable precisely
    when its cyclic momentum is zero.
\end{lemma}
\begin{proof}\leanok
    The finite OBC construction supplies a chain with bond bound
    $d^{\lfloor N/2\rfloor}$. Apply the cyclic block averaging of
    \cite[Theorem~3]{PerezGarcia2007Matrix}; exact invariance makes the
    average equal to the original vector. The converse is cyclicity
    of the matrix trace, including every power of the basic cyclic shift.
\end{proof}

\begin{lemma}[Unused periodic bond coordinates]
    \label{lem:periodic_uniform_bond_padding}
    \lean{MPSTensor.padBond, MPSTensor.mpv_padBond, MPSTensor.mpvState_padBond}
    \leanok
    \uses{def:mps_tensor}
    Let $\widetilde A^i=A^i\oplus0_{K-D}$ be the extension of every site
    matrix by zero from dimension $D$ to a dimension $K\ge D$. Then for
    every $N>0$,
    \begin{align}
        V^{(N)}(\widetilde A) & =V^{(N)}(A).
        \label{eq:periodic_uniform_bond_padding}
    \end{align}
    The assertion excludes the empty word, whose trace is the
    bond dimension.
\end{lemma}
\begin{proof}\leanok
    For matrices $X,Y$ of dimension $D$,
    $(X\oplus0_{K-D})(Y\oplus0_{K-D})=XY\oplus0_{K-D}$. Induction on
    nonempty words gives
    $\widetilde A^{\sigma_0}\cdots\widetilde A^{\sigma_{N-1}}
        =A^{\sigma_0}\cdots A^{\sigma_{N-1}}\oplus0_{K-D}$, so
    \begin{align}
        \tr\bigl(\widetilde A^{\sigma_0}\cdots\widetilde A^{\sigma_{N-1}}\bigr)
         & =\tr\bigl(A^{\sigma_0}\cdots A^{\sigma_{N-1}}\bigr).
    \end{align}
\end{proof}

\begin{lemma}[One common bond size on the exact cyclic ground sector]
    \label{lem:uniform_bond_cyclic_ground_sector}
    \lean{MPSTensor.exists_uniform_bond_on_cyclic_ground_sector}
    \leanok
    \uses{thm:eventual_mps_prime_length_dichotomy,
        lem:periodic_finite_cyclic_representation, lem:periodic_uniform_bond_padding}
    Suppose the permitted ground spaces $\ker H_M$, $M\ge L$, are
    one-dimensional, and a bond-$D$ tensor $B$ generates their nonzero
    ground lines for every $M\ge M_0$. Then the single matrix size
    \[
        K=D+M_0d^{M_0}+1
    \]
    suffices to generate each permitted positive-length ground line
    containing a nonzero exactly cyclic-invariant vector, including
    every such line at $M<M_0$. The representing tensor may depend
    on $M$.
\end{lemma}
\begin{proof}\leanok
    At long lengths extend $B$ by unused bond coordinates. At a short
    length use the preceding finite representation, whose dimension is
    $M d^{\lfloor M/2\rfloor}\le M_0d^{M_0}<K$, and then extend it.
    The nonzero ground vector and uniqueness identify the whole
    ground line.
\end{proof}

\begin{lemma}[Exact finite-prefix criterion for an all-volume family]
    \label{lem:uniform_bond_finite_prefix_criterion}
    \lean{MPSTensor.uniform_bond_family_iff_short_cyclic_ground}
    \leanok
    \uses{lem:uniform_bond_cyclic_ground_sector}
    Under the preceding assumptions, the following are equivalent:
    \begin{enumerate}
        \item There is one positive matrix size $K$ such that, at
              every permitted positive volume $M$, some bond-$K$ tensor
              has a nonzero periodic vector generating $\ker H_M$.
        \item At each permitted short positive volume $M<M_0$, the
              ground space contains a nonzero vector fixed exactly by
              every cyclic translation.
    \end{enumerate}
    An invariant ground line alone does not state the second condition:
    translation can act on a line by a nontrivial root of unity.
    Lemma~\ref{lem:seven_sector_concrete_interaction} shows that the second
    condition does not follow from positive frustration-free uniqueness on
    all permitted volumes $M\ge7$. The relation to the source's volume and
    representation conventions is recorded in
    \cite{gap:spwc10_ground_state_dichotomy_quantifiers}.
\end{lemma}
\begin{proof}\leanok
    The first condition supplies nonzero periodic trace vectors, and
    trace cyclicity proves the second. Conversely use the preceding
    common-size construction at the short volumes and the already
    supplied eventual tensor at the long volumes.
\end{proof}

\subsection{The normal product-state branch}
\label{sec:spwc_product_finite_prefix}

\begin{lemma}[Nonvanishing of every positive product-state volume]
    \label{lem:spwc_bond_one_all_volume_nonvanishing}
    \lean{MPSTensor.mpv_apply_of_bondOne,
        MPSTensor.exists_mpv_ne_zero_of_bondOne,
        MPSTensor.mpvState_ne_zero_of_normal_bondOne}
    \leanok
    For a bond-one tensor,
    \[
        V^{(M)}(B)_\sigma=\prod_{j=1}^M B_{\sigma_j}.
    \]
    If one letter is nonzero, this vector is nonzero at every positive
    length. In particular this holds for every normal bond-one tensor;
    no finite-volume nonvanishing hypothesis is needed.
\end{lemma}
\begin{proof}\leanok
    The unique virtual index leaves one product in the trace expansion.
    Choose the same nonzero letter at every physical site. Its product
    is nonzero. Normality supplies such a letter.
\end{proof}

\begin{lemma}[Short ground vectors in the product-state branch]
    \label{lem:spwc_bond_one_finite_prefix_ground}
    \lean{MPSTensor.exists_cyclic_ground_of_normal_bondOne_local_kernel,
        MPSTensor.uniform_bondOne_ground_family_of_normal_local_kernel}
    \leanok
    \uses{lem:spwc_bond_one_all_volume_nonvanishing, lem:local_mps_kernel_all_volumes,
        lem:periodic_finite_cyclic_representation}
    Suppose the normal tensor in the eventual fixed-state construction
    has bond dimension one. Retain the derived inclusion of its
    length-$L$ local MPS space in $\ker h$. Then its periodic
    vector is a nonzero zero-momentum ground vector at every positive
    $M\geq L$, including the whole finite prefix. If each such
    Hamiltonian has a one-dimensional zero-energy space, the same
    bond-one tensor generates every one of these ground spaces.
    No short representation or zero-momentum premise is supplied.
    This treats the product-state branch only; the general short-volume
    obligation remains in \cite{gap:spwc10_ground_state_dichotomy_quantifiers}.
\end{lemma}
\begin{proof}\leanok
    The preceding lemma gives $V^{(M)}(B)\ne0$ for every $M>0$. Since
    $G_L(B)\subseteq\ker h$, Lemma~\ref{lem:local_mps_kernel_all_volumes}
    gives $h_iV^{(M)}(B)=0$ for every cyclic translate, so
    \begin{align}
        H_MV^{(M)}(B) & =\sum_{i=0}^{M-1}h_iV^{(M)}(B)=0.
    \end{align}
    Cyclicity of the trace gives
    $V^{(M)}(B)_{\sigma(\cdot+s)}=V^{(M)}(B)_\sigma$ for every shift $s$.
    When $\dim\ker H_M=1$, the nonzero vector $V^{(M)}(B)$ spans the
    kernel.
\end{proof}

\begin{lemma}[Product-state completion from one prime-length ground vector]
    \label{lem:spwc_bond_one_prime_ground_family}
    \lean{MPSTensor.exists_uniform_bondOne_ground_family_of_prime_length_bound}
    \leanok
    \uses{lem:spwc_bond_one_finite_prefix_ground,
        lem:prime_small_representation_fixed_tensor}
    Let $h$ be the one fixed positive interaction of range $L$, and
    assume every permitted periodic Hamiltonian has a one-dimensional
    zero-energy space. If one prime $N\geq L+4$ has a nonzero
    bond-one MPS ground representation, then one bond-one tensor
    generates all permitted positive-volume ground spaces, including
    the entire finite prefix. The normal tensor and local-kernel
    inclusion are derived rather than supplied.
\end{lemma}
\begin{proof}\leanok
    Apply the prime-length extraction, retaining its normality
    and local-kernel inclusion. Its positive retained dimension is at
    most the input dimension one, hence equals one. Apply the preceding
    product-state construction at every positive volume. The numerical
    threshold is the general fifth-degree threshold at $D=1$:
    $q(1)=(1^2+1)^2=4$ and $3(1-1)(q(1)+1)=0$, so
    $\max\{q(1),3(1-1)(q(1)+1)\}=4$.
\end{proof}

\subsection{Fifth-root Weyl word ingredients}
\label{sec:short_volume_weyl_words}

\begin{lemma}[Degree and reachable Weyl charges]
    \label{lem:short_volume_weyl_charges}
    \lean{MPSTensor.ShortMomentum.countCharge,
        MPSTensor.ShortMomentum.wordCharge,
        MPSTensor.ShortMomentum.residues_eq_of_degree_charge,
        MPSTensor.ShortMomentum.wordCharge_ne_zero_of_length_seven,
        MPSTensor.ShortMomentum.exists_word_seven_of_charge,
        MPSTensor.ShortMomentum.exists_word_eight_of_charge,
        MPSTensor.ShortMomentum.exists_word_of_eight_le_of_charge,
        MPSTensor.ShortMomentum.exists_central_counts_of_eight_le}
    \leanok
    Give three letters the charges $(1,0)$, $(0,1)$ and $(-1,-1)$
    in $(\mathbb Z/5\mathbb Z)^2$. At fixed degree, the total charge
    determines all three letter-count residues. At degree seven the
    reached charges are exactly the nonzero charges; at every degree
    at least eight all charges occur. Every degree at least eight admits
    central counts.
\end{lemma}
\begin{proof}\leanok
    The two charges and the total degree determine three times the
    third count, and three has inverse two modulo five. The least
    residues at degree seven sum to two, seven or twelve; the last
    case is exactly the central residue triple $(4,4,4)$.
    At degree eight their sum is three or eight. Add a multiple of five
    to one count to reach the desired degree. Every integer at least
    eight is a nonnegative combination of three and five.
\end{proof}

\begin{lemma}[Exact charged matrix products]
    \label{lem:short_volume_weyl_products}
    \lean{MPSTensor.ShortMomentum.chargedWeyl,
        MPSTensor.ShortMomentum.chargedWeyl_mul,
        MPSTensor.ShortMomentum.threeLetterTensor,
        MPSTensor.ShortMomentum.evalWord_threeLetterTensor}
    \leanok
    \uses{lem:short_volume_weyl_charges}
    Fix a fifth root of unity $q$ and put
    $W(u,v)_{ij}=q^{ui}\mathbf1_{j=i-v}$. Then
    \[
        W(a,b)W(c,d)=q^{-cb}W(a+c,b+d).
    \]
    The three-letter tensor
    $(W(1,0),W(0,1),q^{-1}W(-1,-1))$ evaluates every word to a
    nonzero scalar multiple of the matrix indexed by its charge.
\end{lemma}
\begin{proof}\leanok
    Each matrix row has one retained intermediate index.
    Multiplying its two character values gives the displayed cocycle.
    Induction on physical words records the accumulated exponent.
\end{proof}

\begin{lemma}[Full span at degree eight and vanishing trace at degree seven]
    \label{lem:short_volume_weyl_normality}
    \lean{MPSTensor.ShortMomentum.sum_rootCharacter_smul_chargedWeyl,
        MPSTensor.ShortMomentum.span_chargedWeyl_eq_top,
        MPSTensor.ShortMomentum.isNBlkInjective_threeLetterTensor_eight,
        MPSTensor.ShortMomentum.isNBlkInjective_threeLetterTensor_of_eight_le,
        MPSTensor.ShortMomentum.isNormal_threeLetterTensor,
        MPSTensor.ShortMomentum.trace_evalWord_threeLetterTensor_seven,
        MPSTensor.ShortMomentum.not_isNBlkInjective_threeLetterTensor_seven,
        MPSTensor.ShortMomentum.exists_trace_evalWord_ne_zero_of_eight_le}
    \leanok
    \uses{lem:short_volume_weyl_charges, lem:short_volume_weyl_products}
    Let $q$ be a primitive fifth root of unity.
    The three-letter tensor is injective at every degree at least eight and
    therefore normal. Every degree-seven word is traceless, so it is
    not injective at degree seven. At every degree at least eight its
    periodic trace vector is nonzero. These matrix facts do not classify
    the ground spaces of any proposed fixed local interaction.
\end{lemma}
\begin{proof}\leanok
    The finite Fourier identity
    \[
        \sum_u q^{-ui}W(u,i-j)=5E_{ij}
    \]
    gives all matrix units. Reachability and nonzero word scalars put
    them in the degree-eight span. Noncentral Weyl matrices are
    traceless, and every degree-seven word has noncentral charge.
    At larger degrees, central counts give a nonzero scalar times
    the identity and hence a nonzero trace coefficient.
\end{proof}

\begin{lemma}[Homogeneous Weyl-letter relations]
    \label{lem:short_volume_weyl_letter_relations}
    \lean{MPSTensor.ShortMomentum.chargedWeyl_commute,
        MPSTensor.ShortMomentum.threeLetterTensor_commute,
        MPSTensor.ShortMomentum.threeLetterTensor_cyclic_commute,
        MPSTensor.ShortMomentum.threeLetterTensor_pow_five,
        MPSTensor.ShortMomentum.evalWord_five_letter_transfer}
    \leanok
    \uses{lem:short_volume_weyl_products}
    The letters $X,Y,Z$ satisfy
    \[
        XY=qYX,\qquad YZ=qZY,\qquad ZX=qXZ,
        \qquad X^5=Y^5=Z^5=1.
    \]
    Thus replacing five consecutive equal letters by five copies of
    another letter preserves the matrix word in every left
    and right context. Completeness of these relations for periodic
    local constraints is not asserted.
\end{lemma}
\begin{proof}\leanok
    The charged multiplication formula gives the commutator phase as
    the symplectic pairing. Each cyclic adjacent pair has pairing one.
    A five-letter constant word has zero charge and zero accumulated
    exponent. Concatenation preserves the equality in arbitrary contexts.
\end{proof}

\begin{lemma}[Exact local charged boundary spaces]
    \label{lem:short_volume_weyl_local_spaces}
    \lean{MPSTensor.ShortMomentum.chargedWeylDual,
        MPSTensor.ShortMomentum.trace_chargedWeyl_mul_dual,
        MPSTensor.ShortMomentum.chargeBoundary,
        MPSTensor.ShortMomentum.mem_groundSpace_threeLetterTensor_iff,
        MPSTensor.ShortMomentum.groundSpace_threeLetterTensor_same_charge}
    \leanok
    \uses{lem:short_volume_weyl_products}
    Set $D_{u,v}=q^{-uv}W(-u,-v)/5$. These matrices satisfy
    $\operatorname{tr}(W(p)D_q)=\mathbf1_{p=q}$.
    At every length, the local boundary MPS space consists
    exactly of the vectors whose coefficients are the accumulated
    word phase times an arbitrary function of the total charge.
    This is a local-space characterization; no global periodic
    intersection conclusion is asserted.
\end{lemma}
\begin{proof}\leanok
    Multiply the charged matrices and use the vanishing trace
    at every nonzero charge. The two central phases cancel.
    Given a function $g$ of the charge, the boundary matrix
    $\sum_p g(p)D_p$ represents the asserted coefficient family.
    Conversely every boundary matrix defines such a function by
    its trace pairing with the charged matrices.
\end{proof}

\begin{lemma}[Charge comparison across a cyclic cut]
    \label{lem:short_volume_weyl_cyclic_phase}
    \lean{MPSTensor.ShortMomentum.letterPairing,
        MPSTensor.ShortMomentum.wordPhase_append,
        MPSTensor.ShortMomentum.wordPhase_cons_sub_snoc,
        MPSTensor.ShortMomentum.exists_charge_comparison_of_cyclic_agreement}
    \leanok
    \uses{lem:short_volume_weyl_local_spaces}
    Moving a first letter of charge $(a,b)$ to the end changes the
    scalar exponent by $av-ub$, where $(u,v)$ is the total charge.
    If two local boundary vectors agree across that cyclic cut,
    their charge functions satisfy the corresponding character relation.
    The relation is derived from their coefficients.
\end{lemma}
\begin{proof}\leanok
    Concatenation gives the same cocycle as matrix multiplication.
    The first-letter and last-letter expressions differ by the displayed
    pairing. Cancel the nonzero character of the remaining word in the
    two local-space parametrizations.
\end{proof}

\begin{lemma}[Zero main-sector periodic space at length seven]
    \label{lem:short_volume_weyl_seven_periodic_zero}
    \lean{MPSTensor.ShortMomentum.exists_distinct_seven_cyclic_phases,
        MPSTensor.ShortMomentum.eq_zero_of_seven_cyclic_agreement,
        MPSTensor.ShortMomentum.eq_zero_of_seven_cyclic_local_membership,
        MPSTensor.ShortMomentum.chainGroundSpace_threeLetterTensor_seven_eq_bot}
    \leanok
    \uses{lem:short_volume_weyl_cyclic_phase, lem:short_volume_weyl_charges}
    The seven-site periodic parent space formed from the
    seven-site local boundary space of the three-letter Weyl tensor
    is zero. More precisely, two seven-site boundary vectors agreeing
    across the cyclic cut must both vanish. This is an exceptional
    short-volume obstruction; long-volume uniqueness remains separate.
\end{lemma}
\begin{proof}\leanok
    Each of the 24 noncentral charges has two explicit degree-seven words
    whose first letters have different pairings with that charge.
    The preceding cyclic comparison would give two distinct character
    values times the same charge coefficient. Thus that coefficient
    vanishes. The central charge never occurs at this degree, so every
    vector coefficient is zero. The windows starting at zero
    and one supply the two required local boundary vectors.
\end{proof}

\begin{lemma}[All-degree cyclic closure of boundary spaces]
    \label{lem:short_volume_weyl_global_boundary_closure}
    \lean{MPSTensor.ShortMomentum.exists_six_config_of_charge,
        MPSTensor.ShortMomentum.exists_distinct_eight_cyclic_phases,
        MPSTensor.ShortMomentum.exists_distinct_nine_cyclic_phases,
        MPSTensor.ShortMomentum.exists_distinct_cyclic_phases,
        MPSTensor.ShortMomentum.exists_smul_mpv_of_cyclic_agreement,
        MPSTensor.ShortMomentum.mem_mpvSubmodule_of_cyclic_local_membership}
    \leanok
    \uses{lem:short_volume_weyl_cyclic_phase,
        lem:short_volume_weyl_seven_periodic_zero}
    At every degree at least seven, two global boundary vectors
    agreeing across a cyclic cut close to the periodic trace line.
    This conclusion presupposes only their global boundary-space
    membership and their coefficient agreement. It does not state that
    seven-site local constraints already give global membership.
\end{lemma}
\begin{proof}\leanok
    Explicit word pairs detect distinct cyclic phases at every noncentral
    charge in degrees seven, eight and nine. Appending one copy of each
    letter preserves their charge and extends the witnesses by three
    sites, covering every larger degree. Hence all noncentral charge
    coefficients vanish. The remaining central coefficient is the
    periodic trace coefficient divided by five.
\end{proof}

\begin{lemma}[Long-window open intersection]
    \label{lem:short_volume_weyl_long_intersection}
    \lean{MPSTensor.ShortMomentum.groundSpace_intersection_threeLetterTensor_of_eight_le}
    \leanok
    \uses{lem:short_volume_weyl_normality, lem:short_volume_weyl_letter_relations}
    If the common middle of two overlapping boundary windows has at
    least eight sites, their left and right constraints imply
    membership in the next boundary space. No intersection identity
    is supplied. The steps producing eight and nine sites remain separate.
\end{lemma}
\begin{proof}\leanok
    Fullness of the middle word span makes trace pairing
    injective and gives the required intertwining boundary matrices.
    The fourth power of the first physical letter is its inverse.
    Apply the standard one-sided reconstruction with this inverse and
    zero coefficients on the other two physical letters.
\end{proof}

\begin{lemma}[Endpoint charge coefficients on an overlap]
    \label{lem:short_volume_weyl_overlap_charge_coefficients}
    \lean{MPSTensor.ShortMomentum.exists_overlap_charge_coefficients}
    \leanok
    \uses{lem:short_volume_weyl_cyclic_phase, lem:short_volume_weyl_local_spaces}
    Suppose every restriction obtained by fixing the first or last site
    of a vector lies in the corresponding boundary MPS space.
    There are functions of the total charge for each fixed endpoint
    whose product with the full-word phase equals each physical
    coefficient. Their equality on overlapping interior words is not
    supplied as a premise. Values at unreachable endpoint charges
    may remain arbitrary.
\end{lemma}
\begin{proof}\leanok
    Apply the exact local-space characterization to each endpoint
    restriction. Subtract that endpoint's charge, and cancel its scalar
    exponent together with the concatenation cocycle. This gives
    the asserted normalization by the full-word phase on both sides.
\end{proof}

\begin{lemma}[Charge bridges in the short overlaps]
    \label{lem:short_weyl_overlap_charge_bridges}
    \lean{MPSTensor.ShortMomentum.exists_six_overlap_charge_bridge,
        MPSTensor.ShortMomentum.exists_seven_overlap_charge_bridge,
        MPSTensor.ShortMomentum.eq_of_six_overlap_charge_edges,
        MPSTensor.ShortMomentum.eq_of_seven_overlap_charge_edges}
    \leanok
    Write $c_0=(1,0)$, $c_1=(0,1)$ and $c_2=(-1,-1)$ in
    $(\mathbb Z/5\mathbb Z)^2$, and fix a total charge $p$.
    For six-letter middle words, call an edge $(i,j)$ allowed when
    $p-c_i-c_j$ lies outside $\{(0,4),(1,1),(4,0)\}$; an endpoint is
    active when $p\ne c_i$. Every active left endpoint and active right
    endpoint are joined by three allowed edges. For seven-letter middle
    words, an edge is allowed when $p-c_i-c_j\ne0$, and every pair of
    endpoints is joined by three allowed edges.
    Consequently, if $F_j=G_i$ on every allowed edge, the equality holds
    on every active pair in the first case and every pair in the second.
\end{lemma}
\begin{proof}
    \leanok
    For each of the twenty-five charges and each pair of relevant endpoints,
    choose letters $k,l$ with allowed edges $(i,k),(l,k),(l,j)$.
    The finite residue identities give all the witnesses.
    Then $F_j=G_l=F_k=G_i$.
\end{proof}

\begin{lemma}[The two short Weyl intersections]
    \label{lem:short_weyl_eight_nine_intersections}
    \lean{MPSTensor.ShortMomentum.groundSpace_intersection_threeLetterTensor_eight,
        MPSTensor.ShortMomentum.groundSpace_intersection_threeLetterTensor_nine}
    \leanok
    \uses{lem:short_weyl_overlap_charge_bridges,
        lem:short_volume_weyl_overlap_charge_coefficients,
        lem:short_volume_weyl_local_spaces,
        lem:short_volume_weyl_global_boundary_closure,
        lem:short_volume_weyl_charges}
    For the three-letter fifth-root Weyl tensor, a vector on eight sites
    whose first and last seven-site coefficient sections belong to the
    respective boundary spaces belongs to the eight-site boundary space.
    Likewise, a vector on nine sites with both eight-site section conditions
    belongs to the nine-site boundary space.
\end{lemma}
\begin{proof}
    \leanok
    Normalize the endpoint charge coefficients by the full-word phase.
    Every middle word equates the two endpoint coefficients at
    its total charge. The charge bridges make these coefficients agree
    with a single function of that charge.
    On eight sites, a first letter is active because its remaining
    seven-site word cannot have zero charge. Choose any active endpoint
    as the reference. On nine sites every endpoint can serve as reference.
    Multiplying the resulting common charge function by the word phase
    gives the full boundary-space representation.
\end{proof}

\begin{lemma}[Main-sector kernels of the seven-site Weyl interaction]
    \label{lem:short_volume_weyl_periodic_kernel}
    \lean{MPSTensor.ShortMomentum.groundSpace_restriction_intersection_threeLetterTensor,
        MPSTensor.ShortMomentum.mem_groundSpace_of_seven_contiguous_constraints,
        MPSTensor.ShortMomentum.chainGroundSpace_seven_le_groundSpace_threeLetterTensor,
        MPSTensor.ShortMomentum.chainGroundSpace_seven_threeLetterTensor_eq_mpvSubmodule,
        MPSTensor.ShortMomentum.ker_parentHamiltonianES_threeLetterTensor_eq_span,
        MPSTensor.ShortMomentum.ker_parentHamiltonianES_threeLetterTensor_seven_eq_bot,
        MPSTensor.ShortMomentum.mpvState_threeLetterTensor_ne_zero}
    \leanok
    \uses{lem:short_volume_weyl_normality,
        lem:short_volume_weyl_seven_periodic_zero,
        lem:short_volume_weyl_global_boundary_closure,
        lem:short_volume_weyl_long_intersection,
        lem:short_weyl_eight_nine_intersections}
    Let the local interaction be the orthogonal projection onto
    the complement of the seven-site boundary space of the three-letter
    Weyl tensor. Its periodic Hamiltonian has zero kernel at
    length seven. At every length $N\ge8$, its kernel is exactly the
    line generated by its periodic trace vector, which is nonzero.
    No open-chain intersection or global-space premise is supplied.
    This statement concerns the main physical sector alone.
\end{lemma}
\begin{proof}\leanok
    The two short overlap calculations and the full-span reconstruction
    give every one-step boundary-space intersection from range seven.
    Contiguous-window induction yields global boundary membership.
    Cyclic translation preserves the local constraints, so the
    translated vector also lies in the global boundary space. Apply
    the cyclic closure theorem. Finally transport the existing
    equality between cyclic constraints and the parent-Hamiltonian kernel
    to the Euclidean realization. The explicit central words prove
    nonvanishing at every length at least eight.
\end{proof}

\subsection{A fixed interaction with an exceptional short ground line}

The following statements concern the permitted volumes $M\ge7$ and the
ordinary site-independent periodic trace convention. They distinguish an
all-volume representation from the eventual fixed-tensor conclusion in
\cite[Theorem~5]{Sanz2010Quantum}. They make no assertion about a convention
for interactions on fewer than seven sites or about general site-dependent
matrix product representations.

\begin{definition}[The seven-site two-sector interaction]
    \label{def:seven_sector_interaction}
    \lean{MPSTensor.ShortMomentum.sevenSectorInteraction}
    \leanok
    \uses{lem:short_volume_weyl_local_spaces,def:pvbs_pure_window_interaction}
    Split the five physical letters into three Weyl letters and two particle
    letters. Let $E_R$ and $F_R$ be their word embeddings. For parameters
    $\zeta,q\in\mathbb C$, define $h$ to be the orthogonal projection onto
    the complement of
    $E_7\mathcal G_7(\zeta)+F_7\operatorname{span}\{\Psi_7(q)\}$.
    Here $\mathcal G_7(\zeta)$ is the three-letter Weyl boundary space.
\end{definition}

\begin{lemma}[Positivity of the seven-site interaction]
    \label{lem:seven_sector_interaction_positive}
    \lean{MPSTensor.ShortMomentum.sevenSectorInteraction_isPositive}
    \leanok
    \uses{def:seven_sector_interaction}
    The interaction $h$ is positive for all $\zeta,q\in\mathbb C$.
\end{lemma}
\begin{proof}
    \leanok
    It is an orthogonal projection.
\end{proof}

\begin{lemma}[Exact periodic kernels of the two-sector interaction]
    \label{lem:seven_sector_periodic_kernels}
    \lean{MPSTensor.ShortMomentum.ker_sevenSectorHamiltonian_of_eight_le}
    \lean{MPSTensor.ShortMomentum.ker_sevenSectorHamiltonian_seven}
    \leanok
    \uses{def:seven_sector_interaction,lem:physical_sum_full_kernel,
        lem:short_volume_weyl_periodic_kernel,lem:pvbs_pure_window_large_kernel,
        lem:pvbs_matching_volume_kernel}
    Suppose $\zeta$ is a primitive fifth root of unity, and let
    $H_M=\sum_i h_i$ be the periodic Hamiltonian of $h$.
    If $q\ne0$ and $M\ge8$, then
    $\ker H_M=\operatorname{span}\{E_M\psi_M(\zeta)\}$, where
    $\psi_M(\zeta)$ is the Weyl trace vector.
    If $q^7=1$, then
    $\ker H_7=\operatorname{span}\{F_7\Psi_7(q)\}$.
\end{lemma}
\begin{proof}
    \leanok
    Write $H^{W}_M$ and $H^{P}_M$ for the periodic Hamiltonians of the
    two component interactions. Lemma~\ref{lem:physical_sum_full_kernel}
    gives
    \begin{align}
        \ker H_M & =E_M\ker H^{W}_M+F_M\ker H^{P}_M.
    \end{align}
    The Weyl component satisfies $\ker H^{W}_7=0$ and
    $\ker H^{W}_M=\mathbb C\psi_M(\zeta)$ for $M\ge8$; the particle
    component satisfies $\ker H^{P}_7=\mathbb C\Psi_7(q)$ and
    $\ker H^{P}_M=0$ for $M\ge8$.
\end{proof}

\begin{lemma}[A fixed tensor for every larger volume]
    \label{lem:seven_sector_eventual_fixed_tensor}
    \lean{MPSTensor.ShortMomentum.exists_bond_five_eventual_ground_family_sevenSector}
    \leanok
    \uses{lem:seven_sector_periodic_kernels,lem:mpv_state_extend_physical}
    Under the same primitive-fifth-root and $q\ne0$ assumptions, there is
    one site-independent tensor with physical dimension five and bond
    dimension five whose periodic vector is nonzero and spans $\ker H_M$
    for every $M\ge8$.
\end{lemma}
\begin{proof}
    \leanok
    Extend the three-letter Weyl tensor by zero on the two particle letters.
    Its periodic vectors are the isometric extensions of the nonzero
    Weyl trace vectors, and the kernel identity applies at every larger volume.
\end{proof}

\begin{lemma}[Nonzero ground lines at every permitted volume]
    \label{lem:seven_sector_permitted_ground_lines}
    \lean{MPSTensor.ShortMomentum.exists_nonzero_ground_line_sevenSectorHamiltonian}
    \lean{MPSTensor.ShortMomentum.finrank_ker_sevenSectorHamiltonian}
    \leanok
    \uses{lem:seven_sector_periodic_kernels}
    Suppose $\zeta$ is a primitive fifth root of unity and $q^7=1$.
    For every $M\ge7$, there is a nonzero vector $v_M$ with
    $H_Mv_M=0$ and $\ker H_M=\operatorname{span}\{v_M\}$.
    In particular, $\dim\ker H_M=1$.
\end{lemma}
\begin{proof}
    \leanok
    At seven sites use the embedded particle vector, whose first-excitation
    coefficient is one. At every larger volume use the embedded Weyl trace
    vector. Isometric embedding preserves nonvanishing in both cases.
\end{proof}

\begin{lemma}[No invariant or periodic trace vector at the short volume]
    \label{lem:seven_sector_no_short_trace_ground}
    \lean{MPSTensor.ShortMomentum.eq_zero_of_sevenSectorHamiltonian_fixed}
    \lean{MPSTensor.ShortMomentum.mpvState_eq_zero_of_sevenSectorHamiltonian_eq_zero}
    \leanok
    \uses{lem:seven_sector_periodic_kernels,lem:pvbs_matching_volume_no_fixed_ground_vector,
        lem:periodic_finite_cyclic_representation}
    Suppose additionally $q\ne1$. Every vector in $\ker H_7$ fixed exactly
    by every cyclic translation is zero. Consequently, at every bond
    dimension, an ordinary site-independent periodic trace-MPS vector
    killed by $H_7$ is zero.
\end{lemma}
\begin{proof}
    \leanok
    Restrict a vector in the embedded particle line to its particle
    coefficients. Cyclic translation commutes with this restriction, so an
    invariant vector would give an invariant vector in the particle ground
    line, where the nontrivial phase forces it to vanish. Every ordinary
    periodic trace vector is fixed exactly by translation, by trace cyclicity.
\end{proof}

\begin{lemma}[A fixed frustration-free interaction with a short cyclic phase]
    \label{lem:seven_sector_concrete_interaction}
    \lean{MPSTensor.ShortMomentum.exists_seven_site_interaction_without_short_fixed_ground}
    \leanok
    \uses{lem:seven_sector_interaction_positive,lem:seven_sector_permitted_ground_lines,
        lem:seven_sector_eventual_fixed_tensor,lem:seven_sector_no_short_trace_ground}
    There exists one positive seven-site interaction on five physical letters
    with the following properties. At every $M\ge7$ its periodic
    kernel is one-dimensional and contains a nonzero vector killed by every
    local term. One tensor of bond dimension five spans its nonzero
    periodic ground line at every $M\ge8$. At seven sites, no nonzero ground
    vector is fixed exactly by cyclic translation.
\end{lemma}
\begin{proof}
    \leanok
    Choose $\zeta=\exp(2\pi i/5)$ and $q=\exp(2\pi i/7)$ and use the
    two-sector interaction. The preceding kernel identities give the unique
    nonzero lines. For the positive local terms,
    \begin{align}
        \sum_ih_iv=0 & \quad\Longleftrightarrow\quad h_iv=0\ \text{for every }i,
    \end{align}
    so each ground vector is killed by every local term. The seventh root has nontrivial
    cyclic phase, while physical extension gives the eventual
    bond-five tensor.
\end{proof}
